\section{Final Matched Closest-Baseline Campaign}
\label{sec:final-baseline-campaign}

The consolidation questions establish that the frozen method is empirically viable, compatible with multi-step FedAvg updates, and implementable as a complete bidirectional protocol. The remaining empirical closure task is stricter: compare the final method against the closest relevant alternatives under a matched local-learning protocol, a controlled tuning budget, independent held-out partitions, and the same end-to-end communication accounting. This section therefore does not introduce another mechanism. It asks whether the frozen Event-FedAvg operating point remains competitive after the strongest nearby baselines are given a fair opportunity to tune their own compression parameters.

\subsection{Frozen learning protocol and comparator set}

All methods use synchronous full-participation FedAvg on the strong non-IID Fashion-MNIST federation from Section~\ref{sec:consolidation-q3}. Every client performs
\begin{equation}
    E=5
\end{equation}
local SGD steps with
\begin{equation}
    \eta=0.1,
\end{equation}
the same minibatch size, initialization, data partition, and stochastic minibatch stream for a given seed. The two architectures are the ten-class MLP with $d=25818$ parameters and the compact CNN with $d=14538$ parameters.

The comparator set is intentionally small and difficult:
\begin{enumerate}
    \item \textbf{Event-FedAvg V1}, frozen at the Q3 setting and not retuned in this campaign;
    \item \textbf{Strom-style residual pulses}, the closest historical predecessor identified by Q1;
    \item \textbf{EF-TopK}, a strong residual-conserving sparse FedAvg baseline;
    \item \textbf{Sign-EF}, dense sign compression with a transmitted scale and error feedback;
    \item \textbf{dense FedAvg}.
\end{enumerate}

For Event-FedAvg, the MLP uses
\begin{equation}
    \rho=0.999,\quad \Delta=0.025,\quad
    q_r=0.005(1+r/100)^{-0.1},
\end{equation}
and the CNN uses
\begin{equation}
    \rho=0.999,\quad \Delta=0.025,\quad
    q_r=0.01(1+r/100)^{-0.3}.
\end{equation}
The MLP runs for 150 rounds and the CNN for 80 rounds, matching Q3.

The Strom-style FedAvg comparator receives exactly the same local delta representation as V1,
\begin{equation}
    r_i^{r+1}=r_i^r+\frac{p_i}{\eta}\Delta_i^r.
\end{equation}
If $|r_{ij}^{r+1}|\geq\tau$, it sends one coordinate/sign pulse of magnitude $\tau$, the server applies that pulse, and the residual is updated subtractively,
\begin{equation}
    r_{ij}^{r+1}\leftarrow
    r_{ij}^{r+1}-\tau\sign(r_{ij}^{r+1}).
\end{equation}
Thus the comparison isolates the distinction identified by Q1: residual-conserving tied-resolution pulses versus V1's leaky, full-reset, independently resolved pulses.

\subsection{Matched tuning protocol}

Compression settings are selected only on the development partition with partition seed 2400 and training seed 70707. Event-FedAvg remains frozen. EF-TopK is allowed the retained-fraction grid
\begin{equation}
    \kappa\in\{0.005,0.01,0.025,0.05\},
\end{equation}
while Strom is allowed
\begin{equation}
    \tau\in\{0.00125,0.0025,0.005,0.01,0.02\}.
\end{equation}
Each family selects the setting with the smallest final training objective. Sign-EF and dense FedAvg have no compression grid in this campaign.

This procedure deliberately favors predictive quality rather than selecting a low-traffic baseline. It yields
\begin{equation}
\begin{aligned}
\text{MLP:}\qquad &\kappa_{\rm EF}=0.05,\qquad \tau_{\rm Strom}=0.00125,\\
\text{CNN:}\qquad &\kappa_{\rm EF}=0.05,\qquad \tau_{\rm Strom}=0.005.
\end{aligned}
\end{equation}
These settings are then frozen and evaluated on independent partition seeds
\begin{equation}
    2500,\qquad2600,\qquad2700,
\end{equation}
with independent but method-matched training seeds.

The development sweep also exposes a communication--quality trade-off that matters for interpreting Strom. On the CNN, the quality-selected $\tau=0.005$ Strom setting reaches 82.99\% accuracy on the development partition but requires 357.24~Mbit under conservative end-to-end unicast accounting, versus 81.10\% and 71.24~Mbit for V1. Conversely, $\tau=0.02$ reduces Strom traffic to 77.85~Mbit, close to V1, but drops to 80.75\% accuracy. On the MLP, $\tau=0.02$ is almost exactly traffic matched to V1 (182.27 versus 182.17~Mbit on the development partition) but reaches only 76.73\% accuracy versus 82.90\% for V1. These observations motivate the explicit held-out traffic-matched control below.

\subsection{Symmetric end-to-end communication accounting}

The communication model extends Q2 to every baseline symmetrically. Each method exposes the exact server-update representation induced by its uplink:
\begin{itemize}
    \item Event-FedAvg and Strom replay coordinate/sign pulses;
    \item EF-TopK replays coordinate/float32 values;
    \item Sign-EF replays its sign vector and floating-point scale;
    \item dense FedAvg uses the dense model state.
\end{itemize}
Every nonempty packet receives the same 64-bit packet header. A dense float32 checkpoint is always available. For each synchronization round, the protocol uses the cheaper of exact server-update replay and a dense checkpoint. The conservative unicast metric additionally charges a 32-bit catch-up request from every client every round and sends the selected synchronization representation independently to each of the ten clients. The initial dense model synchronization is charged explicitly.

Therefore the primary communication metric in this section is
\begin{equation}
    B_{\rm total}^{\rm unicast}
    =B_{\uparrow}^{\rm packetized}+B_{\downarrow}^{\rm hybrid,unicast}.
\end{equation}
Logical-broadcast totals are also recorded in the result files, but the unicast value is used for the conservative headline comparison.

\subsection{Held-out quality-selected comparison}

Table~\ref{tab:final-baseline-mlp} gives the three-seed MLP result. Values are mean $\pm$ sample standard deviation.

\begin{table}[ht]
\centering
\small
\caption{Final matched MLP comparison over three held-out partitions. Communication is conservative bidirectional unicast traffic with replay/checkpoint fallback.}
\label{tab:final-baseline-mlp}
\begin{tabular}{lrrrr}
\toprule
Method & Test CE & Acc. [\%] & Worst [\%] & Total [Mbit]\\
\midrule
\textbf{Event-FedAvg} & \textbf{0.4784$\pm$0.0022} & \textbf{83.14$\pm$0.14} & \textbf{56.5$\pm$3.3} & \textbf{183.5$\pm$2.0}\\
EF-TopK 5\% & 0.4950$\pm$0.0019 & 82.45$\pm$0.19 & 51.9$\pm$2.6 & 1010.5$\pm$0.0\\
Sign-EF & 0.5259$\pm$0.0018 & 81.11$\pm$0.25 & 51.4$\pm$3.8 & 435.9$\pm$0.0\\
Strom, $\tau=0.00125$ & 0.5206$\pm$0.0086 & 81.81$\pm$0.10 & 50.3$\pm$3.4 & 1610.9$\pm$6.1\\
Dense FedAvg & 0.4967$\pm$0.0015 & 82.23$\pm$0.15 & 52.2$\pm$1.0 & 2487.0$\pm$0.0\\
\bottomrule
\end{tabular}
\end{table}

The MLP result is a particularly strong closure. V1 has the lowest mean test cross entropy, highest mean accuracy, and highest mean worst-class accuracy of the tested methods. At the same time, its conservative end-to-end traffic is approximately
\begin{equation}
    5.51\times
\end{equation}
lower than quality-selected EF-TopK,
\begin{equation}
    8.78\times
\end{equation}
lower than quality-selected Strom, and
\begin{equation}
    13.55\times
\end{equation}
lower than dense FedAvg.

Table~\ref{tab:final-baseline-cnn} gives the corresponding CNN result.

\begin{table}[ht]
\centering
\small
\caption{Final matched compact-CNN comparison over three held-out partitions.}
\label{tab:final-baseline-cnn}
\begin{tabular}{lrrrr}
\toprule
Method & Test CE & Acc. [\%] & Worst [\%] & Total [Mbit]\\
\midrule
Event-FedAvg & 0.5116$\pm$0.0119 & 81.34$\pm$0.46 & \textbf{38.9$\pm$3.4} & \textbf{72.8$\pm$1.2}\\
\textbf{Strom, $\tau=0.005$} & \textbf{0.4993$\pm$0.0100} & \textbf{82.24$\pm$0.27} & 36.4$\pm$11.5 & 359.5$\pm$4.3\\
EF-TopK 5\% & 0.6165$\pm$0.0018 & 77.06$\pm$0.11 & 23.5$\pm$2.6 & 299.5$\pm$0.0\\
Sign-EF & 0.6694$\pm$0.0022 & 76.07$\pm$0.24 & 16.9$\pm$1.3 & 133.5$\pm$0.0\\
Dense FedAvg & 0.6046$\pm$0.0017 & 77.55$\pm$0.08 & 23.1$\pm$0.6 & 749.1$\pm$0.0\\
\bottomrule
\end{tabular}
\end{table}

Figure~\ref{fig:final-baseline-frontier} places the held-out MLP and CNN results on the same communication--accuracy view and makes the final Pareto interpretation visually explicit.

\FinalBaselineFigure

The CNN result is intentionally not summarized as uniform dominance. The quality-selected Strom setting achieves approximately 0.90 percentage points higher mean accuracy and 0.0123 lower mean CE than V1. It pays for that improvement with
\begin{equation}
    \frac{359.5}{72.8}\approx4.94
\end{equation}
times more conservative end-to-end communication. V1 simultaneously has slightly higher mean worst-class accuracy, though the Strom worst-class statistic has substantial seed variability. Thus the relevant conclusion is a Pareto trade-off rather than a claim that V1 always minimizes predictive error.

\subsection{Held-out traffic-matched nearest-neighbor control}

To test the algorithmic distinction from Q1 more directly, the closest sparse baselines are also evaluated at approximately V1-scale communication. Strom uses $\tau=0.02$ for both architectures. EF-TopK uses 1\% sparsity. These settings were identified from the development sweep before the held-out seeds were inspected.

\begin{table}[ht]
\centering
\small
\caption{Traffic-matched held-out comparison. Values are mean $\pm$ standard deviation over the same three partitions.}
\label{tab:final-baseline-traffic}
\begin{adjustbox}{max width=\textwidth}
\begin{adjustbox}{max width=\textwidth}
\begin{tabular}{llrrrr}
\toprule
Architecture & Method & Test CE & Acc. [\%] & Worst [\%] & Total [Mbit]\\
\midrule
MLP & \textbf{Event-FedAvg} & \textbf{0.4784$\pm$0.0022} & \textbf{83.14$\pm$0.14} & \textbf{56.5$\pm$3.3} & 183.5$\pm$2.0\\
MLP & Strom, $\tau=0.02$ & 0.6035$\pm$0.0367 & 79.51$\pm$0.86 & 43.5$\pm$7.4 & \textbf{181.0$\pm$6.7}\\
MLP & EF-TopK 1\% & 0.5278$\pm$0.0025 & 81.07$\pm$0.26 & 49.4$\pm$1.4 & 209.4$\pm$0.0\\
\midrule
CNN & \textbf{Event-FedAvg} & \textbf{0.5116$\pm$0.0119} & \textbf{81.34$\pm$0.46} & \textbf{38.9$\pm$3.4} & 72.8$\pm$1.2\\
CNN & Strom, $\tau=0.02$ & 0.5592$\pm$0.0485 & 79.88$\pm$1.56 & 19.2$\pm$10.8 & 82.0$\pm$1.7\\
CNN & EF-TopK 1\% & 0.6590$\pm$0.0031 & 75.83$\pm$0.26 & 19.4$\pm$1.6 & \textbf{63.9$\pm$0.0}\\
\bottomrule
\end{tabular}
\end{adjustbox}
\end{table}

This control is the most direct empirical support for the Q1 distinction. On the MLP, Strom is almost perfectly traffic matched to V1, with 181.0 versus 183.5~Mbit, but loses 3.63 percentage points of mean accuracy and increases CE from 0.4784 to 0.6035. On the CNN, traffic-matched Strom actually uses about 12.5\% more communication than V1 but remains 1.46 percentage points lower in mean accuracy and has higher CE. EF-TopK can move slightly below V1's CNN traffic at 1\% sparsity, but its predictive degradation is considerably larger.

The quality-selected and traffic-matched Strom results together therefore expose a clear trade-off.
\begin{takeawaybox}[title={Final nearest-neighbor result}]
Residual-conserving tied pulses can buy slightly better CNN accuracy only by moving far toward the high-traffic side of the frontier. At comparable communication, V1 gives the stronger predictive operating point on both tested architectures.
\end{takeawaybox}
This is precisely the empirical behavior expected if decoupling the communication threshold from the applied model quantum is functionally important rather than merely notational.

\subsection{Interpretation and publication-level verdict}

The campaign supports four conclusions.

First, the Event-FedAvg result is not an artifact of comparison against EF-TopK alone. It remains competitive against the much closer Strom residual-pulse construction identified by Q1.

Second, V1 should not be described as uniformly superior to every Strom operating point. The CNN demonstrates a real quality--communication trade-off: tuned Strom achieves slightly better mean CE and accuracy when allowed nearly five times the end-to-end communication. Preserving this negative qualification strengthens rather than weakens the final claim.

Third, at a matched communication budget the picture is much clearer. V1 is better than the traffic-matched Strom and EF controls on both architectures in mean CE, mean accuracy, and mean worst-class accuracy. This result directly supports the practical value of the lossy full-reset and independently controlled update-resolution design established by Q1.

Fourth, the bidirectional conclusion of Q2 remains essential. The headline ratios above use the complete sparse synchronization protocol. They are not uplink-only compression ratios.

The final empirical baseline gate is therefore
\begin{equation}
    \boxed{\text{matched closest-baseline campaign: PASS}.}
\end{equation}
This pass does not upgrade the literature-level Q1 verdict from qualified to strong novelty. Strom remains a close conceptual predecessor. It does, however, close the main empirical objection that the apparent benefit could be reproduced by a tuned residual-conserving threshold-pulse compressor at the same communication budget.

The remaining core publication task is therefore theoretical rather than another mechanism or baseline search: characterize the frozen lossy stateful compressor mathematically, then extract the manuscript from the living report.
